\documentclass[11pt,a4paper]{article}

\usepackage[T1]{fontenc}
\usepackage[utf8]{inputenc}
\usepackage[french]{babel}
\usepackage[margin=2.5cm]{geometry}
\usepackage{booktabs}
\usepackage{siunitx}
\usepackage[table]{xcolor}
\usepackage{hyperref}
\usepackage{pgfplots}
\pgfplotsset{compat=1.18}
\usepackage{caption}
\usepackage{enumitem}
\usepackage{amsmath}
\usepackage{amssymb}

\sisetup{detect-weight=true, detect-family=true}

\hypersetup{
  colorlinks=true,
  linkcolor=blue!50!black,
  urlcolor=blue!50!black,
  citecolor=blue!50!black,
  pdftitle={Note - Enthalpie relative Al4C3},
  pdfauthor={Pr. John Akapulco}
}

\title{\vspace{-1.5cm}Note technique\\[2pt]
\large Diagramme d'enthalpie relative $\Delta H(P)$ d'Al$_4$C$_3$ (référence Pnma-VII)}
\author{Pr. John Akapulco \\ \small \href{mailto:computchem86@gmail.com}{computchem86@gmail.com}}
\date{23 août 2026}

\begin{document}
\maketitle
\vspace{-0.5cm}

\begin{abstract}
\noindent
Cette note complète le rapport \texttt{rapport\_activite\_Al4C3\_HP\_scan.tex} (scan H(P)
complet, 0--100~GPa, sept polymorphes) en présentant les mêmes résultats sous forme
d'enthalpie \emph{relative} : $\Delta H(P) = H_{\text{phase}}(P) - H_{\text{Pnma-VII}}(P)$,
exprimée en meV/atome, où \textbf{Pnma-VII} -- la phase la plus stable à 0~GPa -- sert de
référence et reste fixée à $\Delta H = 0$ sur toute la gamme de pression. Cette normalisation
rend les écarts de stabilité directement lisibles en meV/atome (l'unité usuelle pour comparer
des polymorphes) et fait ressortir sans ambiguïté les deux croisements de stabilité déjà
identifiés : Pnma-VII~$\to$~Pbam à 25.6~GPa, puis Pbam~$\to$~Pnma à 52.3~GPa.
\end{abstract}

\section{Construction du diagramme}

Pour chaque polymorphe $i$ et chaque pression $P$ de la grille (0 à 100~GPa, pas de 10~GPa),
on calcule
\[
\Delta H_i(P) = H_i(P) - H_{\text{Pnma-VII}}(P), \qquad
\Delta H_i(P)\ \text{en meV/atome} = \frac{\Delta H_i(P)\ [\text{eV/f.u.}]}{7}\times 1000,
\]
avec 7 atomes par formule unitaire (Al$_4$C$_3$ = 4~Al + 3~C). Le choix de Pnma-VII comme
référence n'est pas arbitraire : c'est la phase de plus basse enthalpie à 0~GPa dans les sept
candidats étudiés (voir tableau~\ref{tab:rank0}). Soustraire la même courbe $H_{\text{ref}}(P)$
à chaque phase ne déplace aucun des croisements identifiés dans le rapport précédent -- ces
croisements dépendent uniquement de la différence entre deux courbes, invariante par un
changement de référence commun.

\begin{table}[h]
\centering
\begin{tabular}{clS[table-format=-2.4]}
\toprule
Rang & Polymorphe & {$H$/f.u. à 0~GPa (eV)} \\
\midrule
1 & Pnma-VII (référence) & -43.3535 \\
2 & R-3m               & -43.3294 \\
3 & Pnma-III           & -42.3599 \\
4 & Pbam               & -42.2257 \\
5 & Pnma               & -41.5961 \\
6 & anti-CaFe$_2$O$_4$ & -41.5023 \\
7 & I-43d              & -40.2317 \\
\bottomrule
\end{tabular}
\caption{Classement à 0~GPa (jeu complet, 7/7 convergés). Pnma-VII devance R-3m de seulement
0.024~eV/f.u. ($\approx$3.4~meV/atome) et sert de référence pour le diagramme relatif.}
\label{tab:rank0}
\end{table}

\section{Diagramme d'enthalpie relative}

\begin{figure}[h]
\centering
\begin{tikzpicture}
\begin{axis}[
  width=\textwidth, height=8.5cm,
  xlabel={Pression (GPa)},
  ylabel={$\Delta H$/atome (meV), réf. Pnma-VII},
  legend pos=outer north east,
  legend style={font=\small},
  grid=major,
  extra y ticks={0},
  extra y tick style={grid style={solid, black!40}},
]
\addplot table[x=P_GPa,y=dH_meV_per_atom] {../../hp_curves/hp_rel_Pnma-VII.dat};
\addplot table[x=P_GPa,y=dH_meV_per_atom] {../../hp_curves/hp_rel_R-3m.dat};
\addplot table[x=P_GPa,y=dH_meV_per_atom] {../../hp_curves/hp_rel_Pbam.dat};
\addplot table[x=P_GPa,y=dH_meV_per_atom] {../../hp_curves/hp_rel_Pnma-III.dat};
\addplot table[x=P_GPa,y=dH_meV_per_atom] {../../hp_curves/hp_rel_Pnma.dat};
\addplot table[x=P_GPa,y=dH_meV_per_atom] {../../hp_curves/hp_rel_anti-CaFe2O4.dat};
\addplot table[x=P_GPa,y=dH_meV_per_atom] {../../hp_curves/hp_rel_I-43d.dat};
\legend{Pnma-VII (réf.), R-3m, Pbam, Pnma-III, Pnma, anti-CaFe2O4, I-43d}
\end{axis}
\end{tikzpicture}
\caption{Enthalpie relative $\Delta H(P) = H_{\text{phase}}(P) - H_{\text{Pnma-VII}}(P)$, en
meV/atome. La phase de référence est la droite $\Delta H = 0$ ; toute courbe qui passe sous
zéro devient, à cette pression, plus stable que Pnma-VII. R-3m longe la référence de très près
sur 0--20~GPa ($\Delta H < 6$~meV/atome) sans jamais passer sous zéro. Pbam et Pnma-III
franchissent zéro vers 25--26~GPa ; Pnma devient la courbe la plus basse au-delà de 52~GPa.}
\label{fig:hprel}
\end{figure}

\section{Lecture du diagramme}

\begin{itemize}
  \item \textbf{0--20~GPa} : toutes les courbes restent au-dessus de zéro sauf la référence
  elle-même ; R-3m est la plus proche ($\Delta H \in [3, 6]$~meV/atome), les autres phases sont
  nettement plus hautes (Pbam, Pnma-III $\approx$ 130--160~meV/atome ; I-43d $>$ 400~meV/atome).
  \item \textbf{$\approx$25.6~GPa} : Pbam et Pnma-III croisent $\Delta H = 0$ quasi simultanément
  (Pnma-VII cesse d'être la phase de plus basse enthalpie).
  \item \textbf{25--52~GPa} : Pbam (puis Pnma-III, très proche) devient négatif et le reste,
  signalant sa domination sur Pnma-VII dans cette gamme.
  \item \textbf{$\approx$52.3~GPa} : Pnma croise Pbam/Pnma-III et devient la courbe la plus
  basse pour le reste de la gamme étudiée (jusqu'à 100~GPa).
  \item \textbf{I-43d} reste la courbe la plus haute sur toute la gamme, mais son $\Delta H$
  décroît fortement avec la pression (de $+3.12$~eV/f.u. à 0~GPa à $+132$~meV/atome à
  100~GPa), signe d'une convergence progressive des enthalpies sous forte compression.
\end{itemize}

\appendix
\section{Références}

Données sources : \texttt{hp\_curves/results\_hp.csv} (commit \texttt{470e557}, campagne HP
scan 0--100~GPa). Fichiers de données relatifs générés pour cette note :
\texttt{hp\_curves/hp\_rel\_<polymorphe>.dat} (colonnes \texttt{P\_GPa},
\texttt{dH\_meV\_per\_atom}). Voir \texttt{rapport\_activite\_Al4C3\_HP\_scan.tex} pour la
méthodologie DFT complète et le diagramme en enthalpie absolue.

\end{document}
